\documentclass[11pt,letterpaper]{article}

\usepackage[utf8]{inputenc}
\usepackage[T1]{fontenc}
\usepackage[spanish,es-nodecimaldot]{babel}
\usepackage[letterpaper,margin=2.2cm,headheight=15pt]{geometry}
\usepackage[protrusion=true,expansion=false]{microtype}
\usepackage{xcolor}
\usepackage{booktabs}
\usepackage{tabularx}
\usepackage{longtable}
\usepackage{array}
\usepackage{enumitem}
\usepackage{fancyhdr}
\usepackage{hyperref}
\usepackage{xurl}
\usepackage{listings}

\ifdefined\pdfinfoomitdate\pdfinfoomitdate=1\fi
\ifdefined\pdftrailerid\pdftrailerid{}\fi
\ifdefined\pdfsuppressptexinfo\pdfsuppressptexinfo=15\fi

\definecolor{ugblue}{HTML}{006A8D}
\definecolor{deepgreen}{HTML}{123D3A}
\definecolor{softgray}{HTML}{F1F5F4}
\definecolor{midgray}{HTML}{5D6B6B}
\definecolor{success}{HTML}{157A57}
\definecolor{warning}{HTML}{A45A00}

\hypersetup{
  colorlinks=true,
  linkcolor=ugblue,
  urlcolor=ugblue,
  hypertexnames=false,
  pdftitle={Informe Sprint 2 - Seguridad y parámetros},
  pdfauthor={Emily Dayan Robles Alvarado y Lesly Samay Velasquez Anchundia}
}

\pagestyle{fancy}
\fancyhf{}
\lhead{Informe Sprint 2}
\rhead{Prototipo BI asistido por IA}
\cfoot{\thepage}

\setlength{\parindent}{0pt}
\setlength{\parskip}{0.52em}
\setlist[itemize]{leftmargin=1.5em,itemsep=0.22em,topsep=0.3em}
\setlist[enumerate]{leftmargin=1.7em,itemsep=0.32em,topsep=0.35em}
\renewcommand{\arraystretch}{1.18}
\newcommand{\file}[1]{\texttt{#1}}
\newcommand{\statusok}{\textcolor{success}{\textbf{Completado}}}
\newcommand{\statuspending}{\textcolor{warning}{\textbf{Pendiente de cierre}}}
\newcommand{\notebox}[1]{\begin{center}\fcolorbox{ugblue}{softgray}{\begin{minipage}{0.90\textwidth}\textcolor{ugblue}{\textbf{Nota.}} #1\end{minipage}}\end{center}}

\lstdefinestyle{terminal}{
  basicstyle=\ttfamily\small,
  backgroundcolor=\color{softgray},
  frame=single,
  rulecolor=\color{midgray},
  breaklines=true,
  columns=fullflexible,
  keepspaces=true,
  showstringspaces=false,
  xleftmargin=0.4em,
  xrightmargin=0.4em,
  aboveskip=0.6em,
  belowskip=0.6em
}

\begin{document}
\hypersetup{pageanchor=false}

\begin{titlepage}
  \thispagestyle{empty}
  \vspace*{0.8cm}
  {\color{ugblue}\rule{\textwidth}{2.8pt}}
  \vspace{1.2cm}

  {\Large Documentación técnica del producto\\[0.25em]
  Prototipo BI asistido por inteligencia artificial}

  \vfill

  {\color{deepgreen}\fontsize{29}{35}\selectfont\textbf{Informe del Sprint 2}}\\[0.7em]
  {\LARGE Seguridad, administración y configuración inicial}\\[0.6em]
  {\large Prototipo web de inteligencia de negocios asistido por IA}

  \vspace{1.5cm}

  \begin{tabularx}{\textwidth}{@{}>{\bfseries}l X@{}}
    Responsables & Emily Dayan Robles Alvarado y Lesly Samay Velásquez Anchundia \\
    Código documental & INF-SPR-02 \\
    Versión & 1.0.0 \\
    Periodo de referencia & 4 al 11 de septiembre de 2026 \\
    Rama de trabajo & \file{feature/rbac-y-parametros} \\
    Línea base & \file{develop}; promoción posterior mediante pull request \\
    Repositorio & \href{https://github.com/LFXAS/prototipo-bi-ia}{LFXAS/prototipo-bi-ia} \\
    Estado & Aprobado e integrado mediante PR \#28 \\
    Fecha de cierre & 11 de septiembre de 2026 \\
  \end{tabularx}

  \vfill
  {\color{ugblue}\rule{\textwidth}{2.8pt}}
\end{titlepage}

\hypersetup{pageanchor=true}
\pagenumbering{roman}
\tableofcontents
\newpage
\pagenumbering{arabic}

\section{Resumen ejecutivo}

El Sprint 2 transforma la línea base Docker del Sprint 1 en un incremento administrativo utilizable. Se implementa autenticación con JSON Web Token (JWT), autorización basada en roles y permisos (RBAC), gestión de usuarios, roles, permisos y menús, auditoría de acciones administrativas, parámetros operativos controlados y configuración segura de proveedores LLM.

El incremento no adelanta las capacidades analíticas del anteproyecto. Aún no implementa ETL, datamart, exploración de metadatos, dashboard, reportería, predicción ni generación de propuestas BI. La configuración LLM prepara una integración controlada para un sprint posterior; no transmite datos de negocio ni ejecuta SQL.

\section{Objetivo y alcance}

\subsection{Objetivo}

Proveer una administración segura, comprensible y responsive que permita controlar quién accede al prototipo, qué acciones puede realizar y qué configuración LLM aprobada queda disponible para la fase posterior de Copiloto BI.

\subsection{Alcance entregado}

\begin{tabularx}{\textwidth}{@{}p{0.25\textwidth}X p{0.18\textwidth}@{}}
  \toprule
  Capacidad & Entrega del incremento & Estado \\
  \midrule
  Autenticación & Inicio y cierre de sesión, hash Argon2, JWT de vida limitada y cuenta administradora inicial desde el entorno. & \statusok \\
  RBAC & Usuarios, roles, permisos, asociaciones, autorización FastAPI y menús derivados de permisos. & \statusok \\
  Protección & Cuenta/rol base protegidos, invariantes de recuperación y validaciones de dependencias antes de eliminar. & \statusok \\
  Auditoría & Registro inmutable de inicio de sesión, creación, edición, estado, asociaciones, pruebas y eliminación. & \statusok \\
  UX administrativa & Navegación agrupada, plegable, ocultable, responsive, formularios guiados, tablas de casillas y paginación. & \statusok \\
  Parámetros y LLM & Catálogo sin claves libres, credenciales cloud cifradas desde la web, proveedores Gemini/Qwen/Ollama y prueba limitada desde backend. & \statusok \\
  Ollama local & Perfil Docker opcional, modelo ágil \file{qwen2.5:3b} persistente y prueba de disponibilidad real. & \statusok \\
  \bottomrule
\end{tabularx}

\subsection{Exclusiones}

Quedan fuera de este Sprint: recuperación de contraseña por correo, SSO, multiempresa, microservicios, tematización por empresa, ETL, datamart, KPIs, gráficos, alertas, reportería, forecasting y Copiloto BI funcional. Estas exclusiones son decisiones de alcance, no fallas del incremento.

\section{Desarrollo guiado por especificaciones}

Antes de implementar se organizaron los contratos SDD en \file{docs/specs/}. La especificación principal separa las responsabilidades para evitar pantallas mezcladas o reglas de seguridad implícitas.

\begin{longtable}{@{}p{0.28\textwidth}p{0.36\textwidth}p{0.26\textwidth}@{}}
  \toprule
  Documento & Contrato que define & Evidencia vinculada \\
  \midrule
  \endfirsthead
  \toprule
  Documento & Contrato que define & Evidencia vinculada \\
  \midrule
  \endhead
  \url{SPR-02-01-seguridad-rbac.md} & Modelo RBAC, CRUD, protección administrativa, eliminación física validada y auditoría. & API, migración y pruebas. \\
  \url{SPR-02-02-navegacion-y-experiencia.md} & Navegación responsive, agrupación, plegado, estados de pantalla, accesibilidad y controles legibles. & Prueba visual y frontend. \\
  \url{SPR-02-03-parametros-y-llm.md} & Parámetros controlados, proveedores, credenciales web cifradas y prueba LLM segura. & Adaptadores, cifrado, prueba unitaria y prueba real. \\
  \url{SPR-02-04-identidad-visual-y-shell-bi.md} & Tokens visuales y dirección para módulos futuros sin mostrarlos prematuramente. & CSS y revisión visual. \\
  \url{SPR-02-matriz-correcciones.md} & Trazabilidad de observaciones funcionales y decisiones correctivas. & Revisión, pruebas y bitácora. \\
  \bottomrule
\end{longtable}

\notebox{Una especificación no sustituye la revisión humana. Las autoras aprueban, corrigen o rechazan la propuesta antes de que pase a código, pruebas y pull request.}

\section{Arquitectura del incremento}

\subsection{Flujo de seguridad}

\begin{center}
\fcolorbox{ugblue}{softgray}{\begin{minipage}{0.90\textwidth}
\centering
\textbf{React} $\rightarrow$ \textbf{FastAPI} $\rightarrow$ \textbf{JWT y permiso efectivo}\\[0.4em]
Si autoriza: \textbf{servicio RBAC/parámetros} $\rightarrow$ \textbf{PostgreSQL app}\\[0.4em]
Cada mutación $\rightarrow$ \textbf{evento de auditoría inmutable}
\end{minipage}}
\end{center}

React mejora la orientación de la persona usuaria y oculta acciones no autorizadas, pero FastAPI vuelve a comprobar JWT y permiso en cada endpoint. Así, manipular la interfaz no concede privilegios.

\subsection{Modelo RBAC}

Un usuario puede tener varios roles. Un rol agrupa permisos. Los permisos autorizan endpoints y determinan qué menús son visibles. Los códigos técnicos permanecen estables en backend/migraciones; la interfaz presenta nombres, descripciones y tablas de casillas para evitar que la persona escriba IDs o privilegios manualmente.

Las entidades de aplicación son \file{users}, \file{roles}, \file{permissions}, \file{user\_roles}, \file{role\_permissions}, \file{menus}, \file{menu\_permissions}, \file{audit\_events}, \file{parameters} y \file{llm\_configurations}, bajo el esquema PostgreSQL \file{app}.

\subsection{Controles de recuperación y eliminación}

La cuenta y el rol administrativos iniciales poseen una marca persistida de protección. FastAPI rechaza desactivarlos, eliminarlos o dejarlos sin las capacidades mínimas de recuperación. La eliminación física de recursos no protegidos requiere confirmación, comprueba primero dependencias hijas y responde con conflicto explícito cuando deben resolverse asociaciones. Una eliminación permitida genera auditoría con acción \file{delete}; los eventos de auditoría no se modifican ni eliminan desde la interfaz.

La corrección post-cierre para cuentas temporales permite eliminar un usuario no protegido una vez retiradas sus asociaciones usuario--rol, incluso si fue actor de auditoría. La migración \file{20260911\_05} usa \file{ON DELETE SET NULL} sobre la referencia técnica de actor y conserva una etiqueta histórica segura en el evento. La auditoría permanece disponible, mientras la cuenta temporal deja de existir.

\subsection{Evolución reproducible de base de datos}

El Sprint 2 formaliza la separación entre estructura, catálogo base y datos locales. Alembic aplica cambios de esquema versionados; inmediatamente después, el servicio Compose \file{seed} incorpora de manera idempotente permisos, menús, rol administrador y cuenta inicial protegidos. Así, cada programadora puede actualizar \file{develop} sin copiar un volumen PostgreSQL: su base conserva registros locales no protegidos y recibe el contrato técnico nuevo.

La semilla no transporta auditoría, cuentas temporales, claves, tokens, configuraciones LLM ni datos de negocio. Los datos demostrativos de módulos futuros deberán ser sintéticos, opcionales y aprobados mediante su especificación SDD.

\section{Experiencia de usuario}

La administración se integra en un cascarón responsive. En escritorio se puede contraer o restaurar el panel lateral mediante un chevrón compacto superpuesto al borde. En tableta y móvil se usa el patrón hamburguesa. Los grupos \textbf{Seguridad} y \textbf{Parámetros generales} son plegables incluso cuando contienen la ruta activa; Inicio queda como acceso directo.

Las listas reciben paginación desde FastAPI. Roles, usuarios y menús asignan relaciones mediante tablas de casillas con nombre y descripción. Al editar un usuario, la contraseña es opcional: si queda vacía se conserva la existente. Mensajes de éxito usan confirmación verde accesible; errores y validaciones se traducen a lenguaje comprensible y no muestran estructuras técnicas como \file{[object Object]}.

\section{Parámetros generales}

La pantalla Parámetros está preparada para configuraciones operativas aprobadas, pero no es una tabla libre de clave/valor. Un parámetro sólo podrá aparecer cuando exista un módulo consumidor, nombre humano, descripción, tipo, rango y validación definidos. En Sprint 2 no hay todavía un consumidor de parámetros operativos; por ello la pantalla muestra un estado vacío deliberado.

El primer catálogo se definirá con el módulo que lo necesite, por ejemplo ETL, reportería o analítica. Ese sprint deberá incluir la especificación SDD del valor, sus límites, formulario, auditoría y pruebas. No se guardarán contraseñas, tokens, cadenas de conexión ni valores arbitrarios como parámetros generales.

\section{Configuración LLM}

\subsection{Principio de seguridad}

La configuración conserva nombre, proveedor, URL permitida, modelo, estado y resultado seguro de la prueba. Para Gemini y Qwen Cloud, la persona administradora guarda primero la configuración y usa \textbf{Registrar credencial} o \textbf{Reemplazar credencial}. FastAPI cifra la clave; PostgreSQL conserva únicamente el texto cifrado en \file{app.secrets}, la raíz criptográfica permanece en el volumen Docker \file{secret\_key\_data} y React sólo recibe el estado \textbf{Credencial configurada}. Ollama local no requiere clave.

La migración \file{20260918\_06} vincula cada configuración con, como máximo, un secreto. Cambiar de proveedor invalida la credencial anterior; eliminar una configuración inactiva elimina también su secreto. La auditoría registra alta o reemplazo sin incluir la clave. \file{GEMINI\_API\_KEY} y \file{DASHSCOPE\_API\_KEY} se retiraron de Compose y de la plantilla \file{.env.example}.

\subsection{Alternativa local verificable}

El perfil Compose \file{local-llm} incorpora Ollama de forma opcional y privada dentro de la red de contenedores. El modelo local predeterminado es \file{qwen2.5:3b}, elegido para mantener una experiencia ágil en equipos de desarrollo. \file{qwen3:4b} permanece como alternativa de mayor capacidad y mayor consumo. Se opera desde la raíz del repositorio:

\begin{lstlisting}[style=terminal,language=bash]
make ollama-up
make ollama-pull
make ollama-status
make ollama-down
\end{lstlisting}

El modelo se guarda en el volumen \file{ollama\_models}; cada programadora lo descarga una vez en su máquina. No se versiona con Git, no se publica en GHCR y CI/CD no lo descarga. La URL interna que debe registrarse en la pantalla es \file{http://ollama:11434}. FastAPI consulta \file{/api/tags} y sólo confirma éxito si el servicio responde y el modelo seleccionado está descargado. La verificación real del incremento confirmó esa condición sin enviar datos del negocio y sin alterar la configuración cloud activa.

\section{Verificación y evidencia}

\begin{tabularx}{\textwidth}{@{}p{0.30\textwidth}p{0.20\textwidth}X@{}}
  \toprule
  Control & Resultado & Evidencia \\
  \midrule
  Compose & \statusok & Desarrollo, perfil \file{delivery}, perfil \file{local-llm} y entrega GHCR validados con configuración Compose. \\
  Backend & \statusok & Ruff, Mypy y 13 pruebas Pytest Docker; se cubren cifrado, manipulación del texto cifrado, credencial web, validación de URL y modelo Ollama disponible/no descargado. \\
  Frontend & \statusok & ESLint, 4 pruebas Vitest y construcción Vite dentro de Docker; se cubre el formulario de credencial de un solo uso. \\
  LLM local & \statusok & \file{qwen2.5:3b} predeterminado; prueba FastAPI verifica servicio y modelo interno. \\
  Documentación & \statusok & Bitácora, informe Sprint 1, informe Sprint 2 y manual técnico compilados reproduciblemente con LaTeX Docker. \\
  Seguridad operativa & \statusok & Prueba integrada con credencial temporal: cero coincidencias en texto plano, reutilización tras reiniciar y cero secretos huérfanos después de eliminar. \\
  \bottomrule
\end{tabularx}

La comprobación integral se ejecuta con:

\begin{lstlisting}[style=terminal,language=bash]
make verify
\end{lstlisting}

El comando no descarga modelos LLM: valida solamente que el perfil opcional sea una configuración Compose válida. Esta decisión evita que CI consuma almacenamiento, tiempo y recursos de cómputo en cada ejecución.

\section{Riesgos y controles}

\begin{tabularx}{\textwidth}{@{}p{0.27\textwidth}p{0.34\textwidth}X@{}}
  \toprule
  Riesgo & Control aplicado & Resultado esperado \\
  \midrule
  Pérdida de acceso administrativo & Marca protegida e invariantes verificadas en backend. & Siempre existe una vía de recuperación. \\
  Exposición de secretos & Cifrado autenticado, raíz separada, respuesta segura y auditoría sin claves/JWT. & PostgreSQL sólo recibe texto cifrado y ninguna API key vuelve a React. \\
  Borrado inconsistente & Confirmación, validación de hijos y respuesta 409; los eventos de auditoría conservan etiqueta histórica si se borra un actor temporal. & No hay borrado en cascada silencioso ni pérdida de trazabilidad. \\
  Cuota cloud agotada & Ollama local opcional y modelo persistente. & Prueba LLM sin depender de cuota. \\
  Uso excesivo de recursos & Perfil no automático, modelo explícito, límites de contexto y una sola carga. & El entorno normal no se hace más pesado. \\
  Confusión entre módulos & Especificaciones separadas y estado aislado por ruta. & Formularios y mensajes corresponden al recurso activo. \\
  \bottomrule
\end{tabularx}

\section{Cierre y siguientes pasos}

El incremento administrativo quedó construido, verificado y fusionado en \file{develop} mediante el PR \#28. Su promoción posterior siguió el flujo protegido hacia \file{main}; por ello, el cierre documenta un resultado histórico y no instrucciones todavía pendientes.

Los siguientes sprints deben conservar este contrato: primero especificación SDD aprobada; después migraciones, API, interfaz, pruebas y documentación. El próximo módulo funcional recomendado es la introspección determinística de AdventureWorks y sus metadatos, no la generación automática de SQL por IA.

\section{Control documental}

\begin{tabularx}{\textwidth}{@{}>{\bfseries}p{0.28\textwidth}X@{}}
  \toprule
  Campo & Valor \\
  \midrule
  Fuente editable & \file{docs/sprints/sprint-02-seguridad-y-parametros.tex} \\
  Compilación & \file{make sprint2-report} o \file{make docs} \\
  Documentos relacionados & Especificaciones SPR-02, ADR 0002, bitácora y manual técnico. \\
  Revisión & Debe actualizarse si cambia la arquitectura, alcance, controles, pruebas o evidencia del Sprint 2. \\
  \bottomrule
\end{tabularx}

\subsection*{Historial de cambios}

\begin{tabularx}{\textwidth}{@{}l l X@{}}
  \toprule
  Versión & Fecha & Cambio \\
  \midrule
  0.1.0 & 11-09-2026 & Primera versión: documentación del incremento administrativo, RBAC, UX, parámetros, proveedores LLM y alternativa local Ollama. \\
  0.2.0 & 18-09-2026 & Corrección final: registro web de credenciales cloud, almacenamiento cifrado, volumen de clave maestra, migración, auditoría y verificación integrada. \\
  1.0.0 & 24-09-2026 & Normalización profesional del informe y actualización del estado histórico de cierre. \\
  \bottomrule
\end{tabularx}

\end{document}
